\documentclass[10pt,a4paper]{amsart}
\usepackage{amsmath,amssymb,amsthm,graphicx}
\usepackage{verbatim}
\newenvironment{reference}{\comment}{\endcomment}
%\newenvironment{reference}{}{}
\newenvironment{slogan}{\comment}{\endcomment}
\newenvironment{history}{\comment}{\endcomment}
% For commutative diagrams we use Xy-pic
\usepackage[all]{xy}

% We use 2cell for 2-commutative diagrams.
\xyoption{2cell}
% Theorem environments.
%
\theoremstyle{plain}
\newtheorem{theorem}[subsection]{Theorem}
\newtheorem{proposition}[subsection]{Proposition}
\newtheorem{lemma}[subsection]{Lemma}

\theoremstyle{definition}
\newtheorem{definition}[subsection]{Definition}
\newtheorem{example}[subsection]{Example}
\newtheorem{exercise}[subsection]{Exercise}
\newtheorem{situation}[subsection]{Situation}

\theoremstyle{remark}
\newtheorem{remark}[subsection]{Remark}
\newtheorem{remarks}[subsection]{Remarks}

\numberwithin{equation}{subsection}

% Macros
%
\def\lim{\mathop{\mathrm{lim}}\nolimits}
\def\colim{\mathop{\mathrm{colim}}\nolimits}
\def\Spec{\mathop{\mathrm{Spec}}}
\def\Hom{\mathop{\mathrm{Hom}}\nolimits}
\def\Ext{\mathop{\mathrm{Ext}}\nolimits}
\def\SheafHom{\mathop{\mathcal{H}\!\mathit{om}}\nolimits}
\def\SheafExt{\mathop{\mathcal{E}\!\mathit{xt}}\nolimits}
\def\Sch{\mathit{Sch}}
\def\Mor{\mathop{\mathrm{Mor}}\nolimits}
\def\Ob{\mathop{\mathrm{Ob}}\nolimits}
\def\Sh{\mathop{\mathit{Sh}}\nolimits}
\def\NL{\mathop{N\!L}\nolimits}
\def\CH{\mathop{\mathrm{CH}}\nolimits}
\def\proetale{{pro\text{-}\acute{e}tale}}
\def\etale{{\acute{e}tale}}
\def\QCoh{\mathit{QCoh}}
\def\Ker{\mathop{\mathrm{Ker}}}
\def\Im{\mathop{\mathrm{Im}}}
\def\Coker{\mathop{\mathrm{Coker}}}
\def\Coim{\mathop{\mathrm{Coim}}}

% Boxtimes
%
\DeclareMathSymbol{\boxtimes}{\mathbin}{AMSa}{"02}

%
% Macros for moduli stacks/spaces
%
\def\QCohstack{\mathcal{QC}\!\mathit{oh}}
\def\Cohstack{\mathcal{C}\!\mathit{oh}}
\def\Spacesstack{\mathcal{S}\!\mathit{paces}}
\def\Quotfunctor{\mathrm{Quot}}
\def\Hilbfunctor{\mathrm{Hilb}}
\def\Curvesstack{\mathcal{C}\!\mathit{urves}}
\def\Polarizedstack{\mathcal{P}\!\mathit{olarized}}
\def\Complexesstack{\mathcal{C}\!\mathit{omplexes}}
% \Pic is the operator that assigns to X its picard group, usage \Pic(X)
% \Picardstack_{X/B} denotes the Picard stack of X over B
% \Picardfunctor_{X/B} denotes the Picard functor of X over B
\def\Pic{\mathop{\mathrm{Pic}}\nolimits}
\def\Picardstack{\mathcal{P}\!\mathit{ic}}
\def\Picardfunctor{\mathrm{Pic}}
\def\Deformationcategory{\mathcal{D}\!\mathit{ef}}
\usepackage{cleveref}
\crefname{originallemma}{original lemma}{original lemmas}
\Crefname{originallemma}{Original lemma}{Original lemmas}
\makeatletter
\newlabel{decent-spaces-lemma-algebraic-residue-field-extension-closed-point}{{07U5}{source}{}{sourcepointer-1}{}}
\newlabel{decent-spaces-lemma-algebraic-residue-field-extension-closed-point@cref}{{[originallemma][0][]07U5}{[1][0][]source}{}{}{}}
\makeatother
\crefname{originallemma}{original lemma}{original lemmas}
\Crefname{originallemma}{Original lemma}{Original lemmas}
\makeatletter
\newlabel{lemma-vanishing-base-change}{{07V0}{source}{}{sourcepointer-2}{}}
\newlabel{lemma-vanishing-base-change@cref}{{[originallemma][0][]07V0}{[1][0][]source}{}{}{}}
\makeatother
\crefname{originallemma}{original lemma}{original lemmas}
\Crefname{originallemma}{Original lemma}{Original lemmas}
\makeatletter
\newlabel{lemma-vanishing-separate-closed}{{07V1}{source}{}{sourcepointer-3}{}}
\newlabel{lemma-vanishing-separate-closed@cref}{{[originallemma][0][]07V1}{[1][0][]source}{}{}{}}
\makeatother
\crefname{originalreference}{original reference}{original references}
\Crefname{originalreference}{Original reference}{Original references}
\makeatletter
\newlabel{situation-vanishing}{{07UY}{source}{}{sourcepointer-4}{}}
\newlabel{situation-vanishing@cref}{{[originalreference][0][]07UY}{[1][0][]source}{}{}{}}
\makeatother
\crefname{originallemma}{original lemma}{original lemmas}
\Crefname{originallemma}{Original lemma}{Original lemmas}
\makeatletter
\newlabel{spaces-morphisms-lemma-large-enough}{{0488}{source}{}{sourcepointer-5}{}}
\newlabel{spaces-morphisms-lemma-large-enough@cref}{{[originallemma][0][]0488}{[1][0][]source}{}{}{}}
\makeatother
\crefname{originallemma}{original lemma}{original lemmas}
\Crefname{originallemma}{Original lemma}{Original lemmas}
\makeatletter
\newlabel{spaces-morphisms-lemma-universally-injective-local}{{03MX}{source}{}{sourcepointer-6}{}}
\newlabel{spaces-morphisms-lemma-universally-injective-local@cref}{{[originallemma][0][]03MX}{[1][0][]source}{}{}{}}
\makeatother
\crefname{originalsection}{original section}{original sections}
\Crefname{originalsection}{Original section}{Original sections}
\makeatletter
\newlabel{spaces-morphisms-section-points-fields}{{0485}{source}{}{sourcepointer-7}{}}
\newlabel{spaces-morphisms-section-points-fields@cref}{{[originalsection][0][]0485}{[1][0][]source}{}{}{}}
\makeatother
\title{Universally injective structure morphism in the vanishing situation, by separating the closed points of two rational points}
\author{The Stacks Project Authors}
\date{stacks-project:spaces-cohomology.tex:local-snapshot-20261012; GNU FDL 1.2 or later}
\begin{document}
\maketitle
\noindent\textit{Editorial note.} All delivered mathematical text is the author's own, copied from the pinned Stacks Project chapter file spaces-cohomology.tex; no statement or proof step was rewritten, reordered or removed.\par\medskip
\noindent\textit{Editorial note.} Preamble provenance: the pinned chapter file opens with \textbackslash{}input\{preamble\} and carries no preamble text of its own, so the project's single Stacks-wide preamble staged at sources/stacks-spaces-cohomology/preamble.tex is excerpted verbatim as three spans, lines 8-12 (verbatim package and the project's own reference/slogan/history display environments), lines 14-18 (Xy-pic setup) and lines 154-217 (theorem declarations and mathematical macros). The wrapper class and the amsmath/amssymb/amsthm/graphicx base packages are supplied by the assembler. No macro was invented or redefined.\par\medskip
\noindent\textit{Editorial note.} Same-chapter pointers are located at their \textbackslash{}label lines inside the pinned chapter file; cross-chapter pointers are located in the project's own tag list staged in the snapshot (sources/stacks-chow/tags.txt), at the verbatim lines carrying their canonical tags.\par\medskip
\noindent\textit{Allowed standard prerequisites (editorial).}\begin{itemize}
\item Universal injectivity may be checked after an arbitrary field valued base change. Morphisms of Spaces, label spaces-morphisms-lemma-universally-injective-local (invoked at proof line 3106). It reduces the claim to showing that for every field valued point the base change has at most one point.
\item The vanishing situation is stable under base change by an algebra, and separation of disjoint closed subsets. Same chapter spaces-cohomology.tex, labels lemma-vanishing-base-change (invoked at proof line 3107) and lemma-vanishing-separate-closed (invoked at proof line 3127). The first lets the argument assume that the base ring is a field; the second produces a global function separating two distinct closed points, which forces the two \$k\$-valued points to coincide.
\item Points of an algebraic space over a large algebraically closed extension, and rational points as closed points. Morphisms of Spaces, labels spaces-morphisms-lemma-large-enough (invoked at proof line 3117) and Decent Spaces, label decent-spaces-lemma-algebraic-residue-field-extension-closed-point (invoked at proof line 3124). The first makes \$X(K) \textbackslash{}to |X|\$ surjective for a large enough algebraically closed extension, reducing the claim to singleton sets of rational points; the second makes rational points closed, so that the separating function applies to them.
\end{itemize}\medskip
\ifdefined\hypertarget\hypertarget{sourcepointer-1}{}\fi
\noindent\textit{Source reference.} Original lemma 07U5: sources/stacks-chow/tags.txt, source lines 9960--9960. Navigation pointer for \textbackslash{}ref\{decent-spaces-lemma-algebraic-residue-field-extension-closed-point\} used inside the counted span (line 3124); the argument that consumes it is retained in the span.\par\medskip
\ifdefined\hypertarget\hypertarget{sourcepointer-2}{}\fi
\noindent\textit{Source reference.} Original lemma 07V0: sources/stacks-spaces-cohomology/pinned-original-source.tex, source lines 3041--3041. Navigation pointer for \textbackslash{}ref\{lemma-vanishing-base-change\} used inside the counted span (line 3107); the argument that consumes it is retained in the span.\par\medskip
\ifdefined\hypertarget\hypertarget{sourcepointer-3}{}\fi
\noindent\textit{Source reference.} Original lemma 07V1: sources/stacks-spaces-cohomology/pinned-original-source.tex, source lines 3072--3072. Navigation pointer for \textbackslash{}ref\{lemma-vanishing-separate-closed\} used inside the counted span (line 3127); the argument that consumes it is retained in the span.\par\medskip
\ifdefined\hypertarget\hypertarget{sourcepointer-4}{}\fi
\noindent\textit{Source reference.} Original reference 07UY: sources/stacks-spaces-cohomology/pinned-original-source.tex, source lines 2991--2991. Navigation pointer for \textbackslash{}ref\{situation-vanishing\} used inside the counted span (line 3098); the argument that consumes it is retained in the span.\par\medskip
\ifdefined\hypertarget\hypertarget{sourcepointer-5}{}\fi
\noindent\textit{Source reference.} Original lemma 0488: sources/stacks-chow/tags.txt, source lines 5535--5535. Navigation pointer for \textbackslash{}ref\{spaces-morphisms-lemma-large-enough\} used inside the counted span (line 3118); the argument that consumes it is retained in the span.\par\medskip
\ifdefined\hypertarget\hypertarget{sourcepointer-6}{}\fi
\noindent\textit{Source reference.} Original lemma 03MX: sources/stacks-chow/tags.txt, source lines 4765--4765. Navigation pointer for \textbackslash{}ref\{spaces-morphisms-lemma-universally-injective-local\} used inside the counted span (line 3106); the argument that consumes it is retained in the span.\par\medskip
\ifdefined\hypertarget\hypertarget{sourcepointer-7}{}\fi
\noindent\textit{Source reference.} Original section 0485: sources/stacks-chow/tags.txt, source lines 5532--5532. Navigation pointer for \textbackslash{}ref\{spaces-morphisms-section-points-fields\} used inside the counted span (line 3114); the argument that consumes it is retained in the span.\par\medskip
\begin{lemma}
\label{lemma-vanishing-injective}
In Situation \ref{situation-vanishing} the morphism $p : X \to \Spec(A)$ is
universally injective.
\end{lemma}

\begin{proof}
Let $A \to k$ be a ring homomorphism where $k$ is a field. It suffices to
show that $\Spec(k) \times_{\Spec(A)} X$ has at most one point (see
Morphisms of Spaces, Lemma
\ref{spaces-morphisms-lemma-universally-injective-local}).
Using Lemma \ref{lemma-vanishing-base-change} we may assume that $A$
is a field and we have to show that $|X|$ has at most one point.

\medskip\noindent
Let's think of $X$ as an algebraic space over $\Spec(k)$ and let's
use the notation $X(K)$ to denote $K$-valued points of $X$
for any extension $K/k$, see
Morphisms of Spaces, Section \ref{spaces-morphisms-section-points-fields}.
If $K/k$ is an algebraically closed field extension
of large transcendence degree, then we see that $X(K) \to |X|$
is surjective, see Morphisms of Spaces, Lemma
\ref{spaces-morphisms-lemma-large-enough}. Hence, after replacing $k$
by $K$, we see that it suffices to prove that $X(k)$ is a singleton
(in the case $A = k)$.

\medskip\noindent
Let $x, x' \in X(k)$. By Decent Spaces, Lemma
\ref{decent-spaces-lemma-algebraic-residue-field-extension-closed-point}
we see that $x$ and $x'$ are closed points of $|X|$. Hence $x$ and $x'$
map to distinct points of $\Spec(k)$ if $x \not = x'$ by
Lemma \ref{lemma-vanishing-separate-closed}. We conclude that
$x = x'$ as desired.
\end{proof}

\end{document}
